\scanpage{73}
% archive foliation: 72
\uline{Ex 2} Le ``complexe de DeRham à p.d.'' est défini combinatoirement (sans \unsure{connaître} les p.d.) comme objet bigradué de \deleted{$\unsure{K}$} $Q_*$, avec structure d'algèbre \uline{bigraduée} alternée etc. Donc $F_{C^\bullet}$ le transforme en objet de \unsure{la même esp.} On en peut affirmer encore (sans se donner la ``structure p.d'' sur $F$) que le complexe de DeRham associé à $F$ est \unsure{naturel} (\unsure{relat. aux hom, iso,}\ \unsure{compat.}). Le fait que \inserted{\uline{le bicomplexe}} pour tt degré \unsure{$h$|$n$} donné, \unsure{on} trouve une résolution de $k_*=\Psi_*^0$, flasque (i.e. $\Sigma$-injective) sauf au degré vert. maximum $n$ (où on trouve $\unsure{\Psi_*^n|\Phi_*^n}=\Lambda^n\unsure{\Psi_*|\Phi_*}$) --- donc quasi une résolution $\Sigma$-injective tronquée, montre que la coh de sa transformée par $F$ coïncide avec celle de $C^\bullet$ \deleted{\unsure{non tronqué}} \inserted{jusqu'au degré $n$ \unsure{incl.}}, et même --- \unsure{disons-le mieux}, le complexe est homotope (avec classe d'homotopie donnée) au complexe tronqué supérieur de $C^\bullet$ --- en fait, on a un choix canonique (car $\unsure{\operatorname{DR}_*^{n,\bullet}}\inserted{\unsure{\text{degré }n|\text{degré }\le n}}\longrightarrow C_*^\bullet$ hom \uline{unique}).

